% This file is part of the modular source for Stratified Relational Workspace Theory.
% !TeX root = ../Stratified_Relational_Workspace_Theory.tex
\theorypart{IX}{Formal Workspace Definition}

\topic{31}{Certified Stratified Relational Workspace}

Let
\[
\Theta
=
\left(
\varepsilon_{\mathrm{coh}},
\varepsilon_{\mathrm{align}},
\varepsilon_k,
\underline\eta_c,
\underline\eta_{\mathrm{loop}},
\delta_{\mathrm{node}},
\delta_{\mathrm{null}},
\delta_{\mathrm{spec}},
k,
\varepsilon_{\mathrm{pers}}
\right).
\]
Let
\[
\varsigma
=
(\mathcal D,Y^+,\ell,\tau,\mathcal I,\mathcal Q,\Theta)
\]
be the complete declared scientific specification, where
\(\mathcal Q\subseteq\mathcal N_{\mathrm{pred}}(P)\) is a validated null panel.
A dynamic realization contains a certified stratified relational workspace
when there exists a candidate \(W_t\) satisfying all gates required for the
claimed level under \(\varsigma\). The detailed robustness and identifiability
requirements for varying \(\varsigma\) are developed in Parts~XVII and XVIII.

\subtopic{Structural coherence}
\[
\delta_{\mathrm{coh}}(W_t)
\leq
\varepsilon_{\mathrm{coh}}.
\]

\subtopic{Interventional structural--operational alignment}
\[
\delta_{\mathrm{align}}^{\mathcal I}(W_t)
\leq
\varepsilon_{\mathrm{align}}.
\]

\subtopic{Causal efficacy}
\[
\eta_{\mathrm{causal}}^\nu
\geq
\underline\eta_c.
\]

\subtopic{Loop-mediated reentry, when claimed}
\[
\eta_{\mathrm{loop}}(\tau)
\geq
\underline\eta_{\mathrm{loop}}.
\]

\subtopic{Node-level joint advantage}
\[
\Delta_{\mathrm{joint}}^{\ell}
\geq
\delta_{\mathrm{node}}.
\]

\subtopic{Matched-null transfer dependence}
\[
\Delta_{\mathrm{tr}}^{\min}(\mathcal Q)
\geq
\delta_{\mathrm{null}}.
\]

\subtopic{Causal specificity}
\[
\Delta_{\mathrm{specific}}(r)
\geq
\delta_{\mathrm{spec}}.
\]

\subtopic{Distributedness}
\[
k_{\varepsilon_k}
\geq
k.
\]

\subtopic{Persistence, when claimed}
\[
d_{\mathrm{id}}(W_t,W_{t+\tau})
\leq
\varepsilon_{\mathrm{pers}}.
\]

Symbolically,
\[
\boxed{
\SRW(W_t;\varsigma)
=
\bigwedge_{q\in\mathcal J(\varsigma)}
\operatorname{Gate}_q(W_t;\varsigma).
}
\]
The gate set is claim-relative. A one-step operational workspace need not pass
persistence or reentry gates, while a persistent reentrant workspace must.
No scalar total score is sufficient: failure of a required gate cannot be
hidden by high performance on another gate.

\begin{remark}[Structural, predictive, causal, and temporal claims]
The formal hierarchy now distinguishes four logically independent questions:
\[
\begin{aligned}
&\text{Does a stable shared coordinate exist?}\\
&\text{Does it retain the task-relevant conditional law?}\\
&\text{Does manipulating it change downstream behavior specifically?}\\
&\text{Does it persist or reenter through time?}
\end{aligned}
\]
The preceding structural and causal analysis answers the first three at the level of definitions and conditional
theorems. The following parts develop the fourth.
\end{remark}
